% Questões do ENEM citadas em "Física no ENEM: revisão rápida", reproduzidas dos cadernos oficiais do INEP.
% Gerar: (python-dados) uv run --with pymupdf python ../revisao-enem-fisica/questoes/recortar.py
%        python montar.py ; pdflatex questoes-enem-citadas.tex (duas vezes)
\documentclass[11pt]{article}
% Estilo do material "Física no ENEM: revisão rápida" (pdfLaTeX; MiKTeX ou TeX Live)
\usepackage[T1]{fontenc}
\usepackage[utf8]{inputenc}
\usepackage[brazilian]{babel}
\DeclareUnicodeCharacter{2011}{\mbox{-}}   % hífen não separável
\usepackage[a4paper,top=14mm,bottom=17mm,left=17mm,right=17mm,headheight=14pt,footskip=9mm]{geometry}
\usepackage[sfdefault,lining]{FiraSans}
\usepackage[scaled=1.06]{newtxsf}    % matemática sem serifa, na altura do texto
\usepackage[scaled=0.9]{FiraMono}     % evita fonte bitmap (Type 3) em \texttt e \url
\usepackage{microtype}
\ifdefined\pdfglyphtounicode\pdfglyphtounicode{Delta}{0394}\pdfglyphtounicode{Omega}{03A9}\pdfglyphtounicode{mu}{03BC}\fi   % busca por letras gregas no PDF
\usepackage[table]{xcolor}
\usepackage{graphicx}
\usepackage{tikz}
\usetikzlibrary{arrows.meta,calc,positioning,decorations.pathmorphing,decorations.markings,patterns,shapes.geometric,shapes.misc,fit,backgrounds}
\usepackage{pgfplots}
\pgfplotsset{compat=1.18}
\usepackage[RPvoltages]{circuitikz}
\usepackage[most]{tcolorbox}
\usepackage{fontawesome5}
\usepackage{enumitem}
\usepackage{multicol}
\usepackage{booktabs}
\usepackage{array}
\usepackage{colortbl}
\usepackage{fancyhdr}
\usepackage{needspace}
\usepackage{amsmath}
\usepackage[hidelinks,unicode,bookmarksopen=true,bookmarksnumbered=false]{hyperref}

% ---------------------------------------------------------------- cores (trio validado para daltonismo: azul, laranja, verde-água)
\definecolor{corA}{HTML}{2A78D6}      % assunto 1: mecânica
\definecolor{corB}{HTML}{EB6834}      % assunto 2: eletricidade
\definecolor{corC}{HTML}{1BAF7A}      % assunto 3: ondas
\definecolor{tinta}{HTML}{16202A}     % texto
\definecolor{cinza}{HTML}{4F5A64}     % texto secundário
\definecolor{cinzaclaro}{HTML}{EEF1F4}
\definecolor{alerta}{HTML}{C0392B}
\definecolor{ouro}{HTML}{B7791F}
\colorlet{tema}{corA}                 % cor do capítulo corrente
\colorlet{rodape}{cinza}
\newcommand{\corcap}[1]{\colorlet{tema}{#1}}
\color{tinta}
% matemática acompanha o negrito do texto (LaTeX 2020-10 ou mais novo)
\AddToHook{bfseries}{\ifmmode\else\boldmath\fi}
\AddToHook{mdseries}{\ifmmode\else\unboldmath\fi}

% ---------------------------------------------------------------- página: texto alinhado à esquerda, sem hifenização (leitura mais fácil)
\setlength{\parindent}{0pt}
\setlength{\parskip}{5pt plus 1pt minus 1pt}
\linespread{1.08}
\hyphenpenalty=10000 \exhyphenpenalty=10000
\raggedright
\makeatletter\g@addto@macro\@parboxrestore{\raggedright}\makeatother   % minipages e \parbox também
\raggedbottom
\pagestyle{fancy}\fancyhf{}
\renewcommand{\headrulewidth}{0pt}
\newcommand{\nomecap}{}
\fancyfoot[L]{\footnotesize\color{cinza}Física no ENEM · revisão rápida\ifx\nomecap\empty\else\quad\textcolor{rodape}{\rule[-0.2mm]{2.2mm}{2.2mm}}\ \textcolor{rodape}{\bfseries\nomecap}\fi}
\fancyfoot[R]{\footnotesize\color{cinza}\thepage}
\setlist{leftmargin=1.4em,itemsep=2pt,topsep=3pt,parsep=0pt}

% ---------------------------------------------------------------- títulos
\newcounter{assunto}
% \assunto{cor}{Título}{frase curta}{nome curto para rodapé e marcador}
\newcommand{\assunto}[4]{%
  \clearpage\corcap{#1}\colorlet{rodape}{#1}\renewcommand{\nomecap}{#4}\refstepcounter{assunto}%
  \pdfbookmark[1]{#4}{assunto\theassunto}%
  \begin{tcolorbox}[enhanced,colback=tema,colframe=tema,arc=3mm,boxrule=0pt,left=6mm,right=6mm,top=5mm,bottom=5mm,halign=flush left]
    {\color{white}\small\bfseries ASSUNTO \theassunto\ DE 3}\par\vspace{1mm}
    {\color{white}\fontsize{24}{27}\selectfont\bfseries #2\par}\vspace{1.5mm}
    {\color{white}\large #3\par}
  \end{tcolorbox}\vspace{1mm}}
% \titulo{nome do marcador}{Título grande}: abre página de seção geral (fora dos 3 assuntos)
\newcounter{tituloger}
\newcommand{\titulo}[2]{\corcap{tinta}\colorlet{rodape}{cinza}\renewcommand{\nomecap}{#1}\stepcounter{tituloger}%
  \pdfbookmark[1]{#1}{titulo\thetituloger}{\fontsize{24}{27}\selectfont\bfseries #2\par}\nobreak\vspace{1mm}}
\newcounter{marcsec}
\newcommand{\secao}[1]{\par\Needspace{55mm}\vspace{3mm}\stepcounter{marcsec}\pdfbookmark[2]{#1}{sec\themarcsec}{\Large\bfseries\color{tema}#1\par}\nobreak\vspace{1mm}}
\newcommand{\subsecao}[1]{\par\Needspace{62mm}\vspace{1.5mm}{\large\bfseries\color{tinta}#1\par}\nobreak}
\newcommand{\destaque}[1]{{\bfseries\color{tema}#1}}

% ---------------------------------------------------------------- caixas (inquebráveis: uma ideia nunca fica partida entre páginas)
\tcbset{base/.style={enhanced,arc=2.5mm,boxrule=0.8pt,left=4mm,right=4mm,top=3mm,bottom=3mm,halign=flush left,
  fonttitle=\bfseries,before skip=3.2mm,after skip=3.2mm},
  fileira/.style={raster columns=#1,raster equal height,raster column skip=3mm,raster before skip=3.2mm,raster after skip=3.2mm,raster valign=top},
  fileira/.default=3}
% analogia do dia a dia
\newtcolorbox{analogia}[1][Pense assim]{base,colback=tema!8,colframe=tema!55,title={\faLightbulb[regular]\ \ #1},coltitle=tinta,
  attach boxed title to top left={xshift=4mm,yshift=-2.6mm},boxed title style={colback=tema!22,colframe=tema!55,arc=2mm,boxrule=0.8pt},top=5mm}
% fórmula-chave (cartão)
\newtcolorbox{formula}[2][]{base,colback=white,colframe=tema,boxrule=1.2pt,title={\strut #2},colbacktitle=tema,coltitle=white,
  halign=flush center,left=3mm,right=3mm,#1}
% cartão simples (definições, leis de Newton, fenômenos)
\newtcolorbox{cartao}[1][]{base,colback=tema!8,colframe=tema!55,left=3mm,right=3mm,top=2.5mm,bottom=2.5mm,#1}
% cartão com faixa de título
\newtcolorbox{quadro}[2][]{base,colback=white,colframe=tema,title={\strut #2},colbacktitle=tema,coltitle=white,#1}
% questão real do ENEM (resumida)
\newtcolorbox{enemreal}[2][]{base,colback=cinzaclaro,colframe=tinta!70,title={\faGraduationCap\ \ Caiu no ENEM · #2},
  colbacktitle=tinta!85,coltitle=white,#1}
% exemplo do dia a dia
\newtcolorbox{exemplo}[2][]{base,colback=tema!6,colframe=tema!55,title={\faHome\ \ #2},colbacktitle=tema!55,coltitle=tinta,#1}
% pegadinha / erro comum
\newtcolorbox{pegadinha}[1][Pegadinha]{base,colback=alerta!6,colframe=alerta!70,title={\faExclamationTriangle\ \ #1},
  colbacktitle=alerta!80,coltitle=white}
% macete para decorar
\newtcolorbox{macete}[1][Para decorar]{base,colback=ouro!9,colframe=ouro!75,title={\faBrain\ \ #1},
  colbacktitle=ouro!85,coltitle=white}
% questão de treino
\newcounter{treino}
\newtcolorbox{treino}[1]{base,colback=white,colframe=tema,boxrule=1.1pt,
  title={\refstepcounter{treino}Questão \thetreino\ \ ·\ \ #1},colbacktitle=tema,coltitle=white}
% comentário do gabarito
\newtcolorbox{comentario}[2]{base,colback=white,colframe=#1,boxrule=1pt,title={#2},colbacktitle=#1,coltitle=white,
  before skip=2mm,after skip=2mm,top=1.2mm,bottom=1.2mm}
% coluna do resumo de bolso
\newtcolorbox{bolso}[2]{base,colback=white,colframe=#1,boxrule=1.1pt,title={#2},colbacktitle=#1,coltitle=white,
  left=2.5mm,right=2.5mm,top=2mm,bottom=2mm,fontupper=\small}

\newcommand{\passo}[1]{\par\smallskip\textbf{\color{tema}#1}\ }
\newcommand{\resp}[1]{\par\smallskip{\bfseries\color{tema}\faCheckCircle\ Resposta: #1}}
% onde o aluno erra (nota em destaque dentro das caixas de questão)
\newcommand{\erro}[1]{\par\smallskip{\small\textcolor{alerta}{\faExclamationTriangle\ \textbf{Onde o aluno erra:}}\ #1}}
\newcommand{\nota}[1]{\par\smallskip{\small\color{cinza}#1}}
% alternativas A–E
\newlist{alternativas}{enumerate}{1}
\setlist[alternativas]{label=\protect\circulo{\Alph*},leftmargin=2.2em,itemsep=1.5pt,topsep=4pt}
\newcommand{\circulo}[1]{\tikz[baseline=(c.base)]\node[circle,draw=tema,line width=0.7pt,inner sep=1.1pt,font=\footnotesize\bfseries,text=tema](c){#1};}
\newcommand{\leg}[1]{{\footnotesize\color{cinza}#1}}
\newcommand{\numerao}[2]{\centering{\fontsize{26}{28}\selectfont\bfseries\color{tema}\strut #1}\par{\small\color{cinza}#2}}
% selo numerado (liga os pares nas figuras de analogia)
\newcommand{\selo}[1]{\tikz[baseline=(s.base)]\node[circle,fill=tema,text=white,inner sep=1pt,font=\scriptsize\bfseries,minimum size=3.6mm](s){#1};}

% ---------------------------------------------------------------- figuras: estilo único
\tikzset{fio/.style={line width=0.8pt},
  seta/.style={-{Latex[length=2.4mm]},line width=1.4pt},
  cota/.style={{Latex[length=1.6mm]}-{Latex[length=1.6mm]},line width=0.6pt,cinza}}
\ctikzset{bipoles/length=1.0cm,bipoles/thickness=1.5}

% ---------------------------------------------------------------- tabelas da análise (equações e conceitos)
\newcommand{\barra}[1]{\tikz[baseline=0.4ex]{\fill[tema,rounded corners=1pt] (0,0) rectangle (#1*0.36,0.26);\node[anchor=west,font=\small\bfseries,inner sep=2pt] at (#1*0.36,0.13) {#1};}}
\newcommand{\linhaeq}[3]{\rule[-3.2mm]{0pt}{9mm}#1 & #2 & \barra{#3}\\}
\newenvironment{tabelaeq}[2]{% #1 cor, #2 título
  \corcap{#1}\begin{tcolorbox}[base,colback=white,colframe=tema,boxrule=1.1pt,title={#2},colbacktitle=tema,coltitle=white,left=2mm,right=2mm,top=1mm,bottom=0.5mm,before skip=3mm,after skip=2mm]
  \small\renewcommand{\arraystretch}{1.0}\rowcolors{1}{tema!6}{white}
  \begin{tabular}{@{\ }>{\raggedright\arraybackslash}m{5.3cm}>{\raggedright\arraybackslash}m{6.3cm}m{4.1cm}@{}}}
  {\end{tabular}\end{tcolorbox}}

\hypersetup{pdftitle={Física no ENEM: questões citadas na revisão},pdfauthor={Professor Erly Aquino},pdflang={pt-BR},
  pdfsubject={Questões do ENEM 2021 a 2025 (caderno azul) citadas na revisão, por assunto, com gabarito oficial}}
\renewcommand{\nomecap}{}
\fancyfoot[L]{\footnotesize\color{cinza}Física no ENEM · questões citadas na revisão\ifx\nomecap\empty\else\quad\textcolor{rodape}{\rule[-0.2mm]{2.2mm}{2.2mm}}\ \textcolor{rodape}{\bfseries\nomecap}\fi}
\definecolor{capaAzul}{HTML}{0B6FC6}
\definecolor{capaCiano}{HTML}{0A9CCB}
\definecolor{capaAmarelo}{HTML}{DBEE26}
\definecolor{capaCartao}{HTML}{DCE8F3}
\newcommand{\pc}[2]{($(current page.north west)+(#1,-#2)$)}
% cabeçalho de cada questão: \cabq{cor}{ano}{número}{assunto}
\newcommand{\cabq}[4]{\par\vspace{3mm}\noindent
  \begin{tcolorbox}[enhanced,colback=#1,colframe=#1,arc=2mm,boxrule=0pt,left=3mm,right=3mm,top=1.3mm,bottom=1.3mm,nobeforeafter]
    {\color{white}\bfseries ENEM #2\quad·\quad caderno azul\quad·\quad questão #3\hfill\normalfont\small #4}
  \end{tcolorbox}\par\nobreak\vspace{2mm}}
% recorte do caderno oficial: \recorte{arquivo}{largura}
\newcommand{\recorte}[2]{{\centering\fboxsep=2.5mm\fcolorbox{cinza!35}{white}{\includegraphics[width=#2]{recortes/#1}}\par}\nobreak\vspace{1.5mm}}
% gabarito: \gabq{cor}{letra}{1 se foi resolvida na revisão}
\newcommand{\gabq}[3]{\noindent\hfill
  \ifnum#3=1 {\footnotesize\color{cinza}resolvida passo a passo na revisão\quad}\fi
  \tcbox[enhanced,colback=#1!10,colframe=#1,arc=2mm,boxrule=0.6pt,left=2.5mm,right=2.5mm,top=0.8mm,bottom=0.8mm,nobeforeafter,tcbox raise base]
    {\small Gabarito oficial: \textbf{#2}}\par\vspace{2mm}}
\begin{document}
% ---------------------------------------------------------------- capa (identidade do Aulão de Véspera)
\begin{titlepage}
\thispagestyle{empty}
\pdfbookmark[1]{Capa}{capa}
\begin{tikzpicture}[remember picture,overlay]
  \node[inner sep=0pt] at (current page.center) {\includegraphics[width=\paperwidth,height=\paperheight]{../capa/fundo-a4.png}};
  \node[white,font=\fontsize{92}{92}\selectfont\bfseries,inner sep=0pt] at \pc{10.5}{4.75} {AULÃO};
  \node[white,font=\fontsize{25}{28}\selectfont\bfseries,inner sep=0pt] at \pc{10.5}{6.95} {de Véspera};
  \fill[capaCartao,rounded corners=7mm] \pc{3.7}{8.0} rectangle \pc{17.3}{11.5};
  \node[anchor=west,text width=7.6cm,align=left,inner sep=0pt] at \pc{4.5}{9.75} {%
    {\color{capaAzul}\fontsize{27}{29}\selectfont\bfseries QUESTÕES\\ DO ENEM}\\[2.5mm]
    {\color{capaCiano}\fontsize{13}{15}\selectfont\bfseries FÍSICA · CITADAS NA REVISÃO RÁPIDA}};
  \node[inner sep=0pt] at \pc{14.25}{9.75} {\includegraphics[width=4.5cm]{../capa/logo-enem.png}};
  \node[white,font=\fontsize{14}{18}\selectfont,text width=14cm,align=center,inner sep=0pt] at \pc{10.5}{12.55}
    {As questões completas, como caíram na prova, separadas por assunto e com o gabarito oficial};
  \foreach \x/\nome/\ic in {5.75/Mecânica/\faCarSide, 10.5/Eletricidade/\faBolt, 15.25/Ondas/\faWaveSquare}{
    \node[fill=capaAmarelo,text=capaAzul!70!black,rounded corners=4.5mm,minimum width=4.3cm,minimum height=0.95cm,font=\bfseries\large,inner sep=0pt]
      at \pc{\x}{13.95} {\ic\ \ \nome};}
  \fill[capaCartao,rounded corners=7mm] \pc{5.2}{15.65} rectangle \pc{15.8}{19.15};
  \node[fill=capaAmarelo,text=capaAzul!70!black,rounded corners=2.5mm,minimum width=1.25cm,minimum height=1.25cm,font=\LARGE,inner sep=0pt] at \pc{10.5}{15.55} {\faFile*};
  \node[text width=10cm,align=center,inner sep=0pt] at \pc{10.5}{17.6} {%
    {\color{capaCiano}\fontsize{14}{17}\selectfont ENEM 2021 a 2025}\\[1.5mm]
    {\color{capaCiano}\fontsize{27}{29}\selectfont\bfseries Caderno azul}\\[2mm]
    {\color{capaCiano}\fontsize{14}{17}\selectfont 2º dia · Ciências da Natureza}};
  \node[capaAmarelo,font=\fontsize{12}{14}\selectfont\bfseries,inner sep=0pt] at \pc{10.5}{20.15} {P\,R\,O\,F\,E\,S\,S\,O\,R};
  \node[white,font=\fontsize{36}{38}\selectfont\bfseries,inner sep=0pt] at \pc{10.5}{21.25} {Erly Aquino};
  \draw[white,line width=0.6pt,opacity=0.55] \pc{4.0}{22.45} -- \pc{17.0}{22.45};
  \fill[white,rounded corners=6mm] \pc{3.3}{22.85} rectangle \pc{17.7}{26.3};
  \node[inner sep=0pt] at \pc{5.6}{24.575} {\includegraphics[width=3.7cm]{../capa/logo-gre.png}};
  \node[inner sep=0pt] at \pc{10.65}{24.575} {\includegraphics[width=5.6cm]{../capa/logo-see.png}};
  \node[inner sep=0pt] at \pc{15.5}{24.575} {\includegraphics[width=3.5cm]{../capa/logo-ete.png}};
  \node[anchor=west,white,font=\small,inner sep=0pt,text width=9cm,align=left] at \pc{2.3}{28.25}
    {Questões reproduzidas dos cadernos oficiais do INEP\\ Acompanha ``Física no ENEM: revisão rápida''};
\end{tikzpicture}
\end{titlepage}
\setcounter{page}{2}
% ---------------------------------------------------------------- apresentação
\titulo{Como usar}{Como usar esta lista}
Aqui estão, inteiras, as questões do ENEM que a revisão ``Física no ENEM: revisão rápida'' cita ou resolve.
Elas foram recortadas dos cadernos oficiais, sem nenhuma alteração.

\begin{quadro}{O que você precisa saber}
\begin{itemize}
  \item \textbf{Prova:} ENEM, 2º dia, Ciências da Natureza, anos de 2021 a 2025.
  \item \textbf{Cor e numeração:} todas são do \textbf{caderno azul} (caderno 7). Nos cadernos de outras cores a questão é a mesma,
        mas o número e a ordem das alternativas mudam.
  \item \textbf{Gabarito:} o oficial do INEP, logo abaixo de cada questão e reunido na página~\pageref{pg:gab}. Tampe o gabarito e tente antes.
  \item \textbf{Divisão:} mecânica, eletricidade e ondas, e dentro de cada área por assunto.
  \item \textbf{``Resolvida passo a passo na revisão'':} a resolução comentada está no material de revisão.
\end{itemize}
\end{quadro}

\nota{Fonte: provas e gabaritos do ENEM publicados pelo INEP (gov.br/inep, Provas e Gabaritos), sob licença Creative Commons
Atribuição-SemDerivações 3.0. Os textos e as imagens de terceiros dentro das questões pertencem aos autores indicados em cada uma.}

\assunto{corA}{Mecânica}{10 questões citadas na revisão. Movimento, forças e energia.}{Mecânica}
{\small Assuntos: Cinemática (5); Leis de Newton e forças (4); Energia, trabalho e potência (1).\par}\medskip
\Needspace{570pt}
\secao{Cinemática}
\cabq{corA}{2022}{100}{Cinemática}
\recorte{2022-100-1.pdf}{331.4pt}
\gabq{corA}{B}{1}

\Needspace{604pt}
\cabq{corA}{2023}{100}{Cinemática}
\recorte{2023-100-1.pdf}{331.4pt}
\gabq{corA}{C}{1}

\Needspace{256pt}
\cabq{corA}{2023}{124}{Cinemática}
\recorte{2023-124-1.pdf}{483.0pt}
\gabq{corA}{B}{1}

\Needspace{629pt}
\cabq{corA}{2024}{119}{Cinemática}
\recorte{2024-119-1.pdf}{331.9pt}
\gabq{corA}{B}{0}

\Needspace{649pt}
\cabq{corA}{2025}{126}{Cinemática}
\recorte{2025-126-1.pdf}{313.4pt}
\gabq{corA}{D}{0}

\Needspace{259pt}
\secao{Leis de Newton e forças}
\cabq{corA}{2021}{100}{Leis de Newton e forças}
\recorte{2021-100-1.pdf}{331.4pt}
\Needspace{447pt}
\recorte{2021-100-2.pdf}{331.4pt}
\gabq{corA}{B}{0}

\Needspace{649pt}
\cabq{corA}{2023}{109}{Leis de Newton e forças}
\recorte{2023-109-1.pdf}{285.5pt}
\gabq{corA}{D}{0}

\Needspace{649pt}
\cabq{corA}{2023}{114}{Leis de Newton e forças}
\recorte{2023-114-1.pdf}{317.9pt}
\gabq{corA}{A}{0}

\Needspace{649pt}
\cabq{corA}{2025}{108}{Leis de Newton e forças}
\recorte{2025-108-1.pdf}{223.6pt}
\gabq{corA}{C}{1}

\Needspace{370pt}
\secao{Energia, trabalho e potência}
\cabq{corA}{2021}{107}{Energia, trabalho e potência}
\recorte{2021-107-1.pdf}{483.0pt}
\gabq{corA}{C}{1}

\assunto{corB}{Eletricidade}{8 questões citadas na revisão. Circuitos, pilhas, potência e energia elétrica.}{Eletricidade}
{\small Assuntos: Circuitos elétricos e Lei de Ohm (6); Potência e energia elétrica (2).\par}\medskip
\Needspace{574pt}
\secao{Circuitos elétricos e Lei de Ohm}
\cabq{corB}{2022}{96}{Circuitos elétricos e Lei de Ohm}
\recorte{2022-96-1.pdf}{170.6pt}
\gabq{corB}{A}{0}

\Needspace{219pt}
\cabq{corB}{2022}{119}{Circuitos elétricos e Lei de Ohm}
\recorte{2022-119-1.pdf}{483.0pt}
\gabq{corB}{A}{0}

\Needspace{407pt}
\cabq{corB}{2022}{128}{Circuitos elétricos e Lei de Ohm}
\recorte{2022-128-1.pdf}{331.4pt}
\Needspace{162pt}
\recorte{2022-128-2.pdf}{331.4pt}
\gabq{corB}{A}{1}

\Needspace{649pt}
\cabq{corB}{2023}{120}{Circuitos elétricos e Lei de Ohm}
\recorte{2023-120-1.pdf}{219.5pt}
\gabq{corB}{C}{0}

\Needspace{649pt}
\cabq{corB}{2024}{131}{Circuitos elétricos e Lei de Ohm}
\recorte{2024-131-1.pdf}{216.8pt}
\gabq{corB}{B}{0}

\Needspace{649pt}
\cabq{corB}{2025}{130}{Circuitos elétricos e Lei de Ohm}
\recorte{2025-130-1.pdf}{227.7pt}
\gabq{corB}{E}{1}

\Needspace{502pt}
\secao{Potência e energia elétrica}
\cabq{corB}{2021}{128}{Potência e energia elétrica}
\recorte{2021-128-1.pdf}{331.4pt}
\gabq{corB}{D}{1}

\Needspace{649pt}
\cabq{corB}{2025}{94}{Potência e energia elétrica}
\recorte{2025-94-1.pdf}{310.4pt}
\gabq{corB}{D}{1}

\assunto{corC}{Ondas}{11 questões citadas na revisão. Som, luz e fenômenos ondulatórios.}{Ondas}
{\small Assuntos: Acústica (4); Ondas eletromagnéticas e radiações (4); Ondas: propriedades e fenômenos (3).\par}\medskip
\Needspace{298pt}
\secao{Acústica}
\cabq{corC}{2021}{92}{Acústica}
\recorte{2021-92-1.pdf}{331.4pt}
\Needspace{244pt}
\recorte{2021-92-2.pdf}{331.4pt}
\gabq{corC}{D}{1}

\Needspace{328pt}
\cabq{corC}{2023}{111}{Acústica}
\recorte{2023-111-1.pdf}{483.0pt}
\gabq{corC}{C}{0}

\Needspace{649pt}
\cabq{corC}{2023}{135}{Acústica}
\recorte{2023-135-1.pdf}{219.5pt}
\gabq{corC}{D}{1}

\Needspace{375pt}
\cabq{corC}{2024}{99}{Acústica}
\recorte{2024-99-1.pdf}{331.9pt}
\gabq{corC}{D}{0}

\Needspace{483pt}
\secao{Ondas eletromagnéticas e radiações}
\cabq{corC}{2021}{120}{Ondas eletromagnéticas e radiações}
\recorte{2021-120-1.pdf}{483.0pt}
\gabq{corC}{A}{0}

\Needspace{253pt}
\cabq{corC}{2023}{112}{Ondas eletromagnéticas e radiações}
\recorte{2023-112-1.pdf}{483.0pt}
\gabq{corC}{A}{1}

\Needspace{639pt}
\cabq{corC}{2024}{106}{Ondas eletromagnéticas e radiações}
\recorte{2024-106-1.pdf}{331.9pt}
\gabq{corC}{D}{0}

\Needspace{649pt}
\cabq{corC}{2025}{135}{Ondas eletromagnéticas e radiações}
\recorte{2025-135-1.pdf}{223.4pt}
\gabq{corC}{A}{0}

\Needspace{491pt}
\secao{Ondas: propriedades e fenômenos}
\cabq{corC}{2021}{131}{Ondas: propriedades e fenômenos}
\recorte{2021-131-1.pdf}{331.4pt}
\gabq{corC}{B}{1}

\Needspace{358pt}
\cabq{corC}{2023}{128}{Ondas: propriedades e fenômenos}
\recorte{2023-128-1.pdf}{331.4pt}
\gabq{corC}{C}{0}

\Needspace{460pt}
\cabq{corC}{2025}{114}{Ondas: propriedades e fenômenos}
\recorte{2025-114-1.pdf}{338.3pt}
\gabq{corC}{B}{0}

\clearpage\titulo{Anuladas}{Questões anuladas citadas na análise}
{\small O INEP anulou estas duas questões. Elas aparecem na contagem da análise, por isso estão aqui. Não há gabarito.\par}\medskip
\Needspace{356pt}
\cabq{cinza}{2024}{129}{Calorimetria e mudanças de fase}
\recorte{2024-129-1.pdf}{331.9pt}
\gabq{cinza}{Anulado}{0}

\Needspace{478pt}
\cabq{cinza}{2025}{132}{Acústica}
\recorte{2025-132-1.pdf}{338.3pt}
\gabq{cinza}{Anulado}{0}

\clearpage\titulo{Gabarito}{Gabarito oficial (caderno azul)}\label{pg:gab}
{\small Fonte: gabaritos oficiais do INEP para o caderno 7 (azul) do 2º dia de cada ano.\par}\medskip
\begin{tcbraster}[fileira=3,raster column skip=3mm]
\begin{bolso}{corA}{Mecânica}
\centering\renewcommand{\arraystretch}{1.25}\begin{tabular}{@{}l@{\quad}l@{\quad}c@{}}
2021 & questão 100 & \textbf{B}\\
2021 & questão 107 & \textbf{C}\\
2022 & questão 100 & \textbf{B}\\
2023 & questão 100 & \textbf{C}\\
2023 & questão 109 & \textbf{D}\\
2023 & questão 114 & \textbf{A}\\
2023 & questão 124 & \textbf{B}\\
2024 & questão 119 & \textbf{B}\\
2025 & questão 108 & \textbf{C}\\
2025 & questão 126 & \textbf{D}
\end{tabular}
\end{bolso}
\begin{bolso}{corB}{Eletricidade}
\centering\renewcommand{\arraystretch}{1.25}\begin{tabular}{@{}l@{\quad}l@{\quad}c@{}}
2021 & questão 128 & \textbf{D}\\
2022 & questão 96 & \textbf{A}\\
2022 & questão 119 & \textbf{A}\\
2022 & questão 128 & \textbf{A}\\
2023 & questão 120 & \textbf{C}\\
2024 & questão 131 & \textbf{B}\\
2025 & questão 94 & \textbf{D}\\
2025 & questão 130 & \textbf{E}
\end{tabular}
\end{bolso}
\begin{bolso}{corC}{Ondas}
\centering\renewcommand{\arraystretch}{1.25}\begin{tabular}{@{}l@{\quad}l@{\quad}c@{}}
2021 & questão 92 & \textbf{D}\\
2021 & questão 120 & \textbf{A}\\
2021 & questão 131 & \textbf{B}\\
2023 & questão 111 & \textbf{C}\\
2023 & questão 112 & \textbf{A}\\
2023 & questão 128 & \textbf{C}\\
2023 & questão 135 & \textbf{D}\\
2024 & questão 99 & \textbf{D}\\
2024 & questão 106 & \textbf{D}\\
2025 & questão 114 & \textbf{B}\\
2025 & questão 135 & \textbf{A}
\end{tabular}
\end{bolso}
\end{tcbraster}

\end{document}
